\documentclass[unknownkeysallowed,xcolor=table]{beamer}
 
\usepackage[T2A,T1]{fontenc}
\usepackage[utf8]{inputenc}
\usepackage[english,russian]{babel}
\usepackage{amsmath}
\usepackage{listings}
\usepackage{url}
\usepackage{textcomp}
\usepackage{multirow}

% to not copy line numbers
\usepackage{accsupp}
% to not copy line numbers

\setbeamertemplate{navigation symbols}{}

\newcommand{\textapprox}{\raisebox{0.5ex}{\texttildelow}}

\colorlet{mygreen}{green!60!blue}
\colorlet{mymauve}{red!60!blue}
\definecolor{pblue}{rgb}{0.1, 0.2, 0.8}

\newcommand{\rarr}{$\rightarrow$}

\lstset{
      basicstyle=\ttfamily\small,
      commentstyle=\color{mygreen},
      keywordstyle=\color{blue},
      numberstyle=\noncopynumber,
      stringstyle=\color{mymauve},
      numbers=left,
      stepnumber=1,
      columns=fullflexible,
      breaklines=true,
      postbreak=\mbox{\textcolor{red}{\ensuremath{\hookrightarrow}\space}},
      literate={~} {\textapprox}{1},
      language={[11]C++}
}

\makeatletter
\newcommand{\srcmediumsize}{\@setfontsize{\srcmediumsize}{7pt}{7pt}}
\makeatother

\makeatletter
\newcommand{\srcbigsize}{\@setfontsize{\srcbigsize}{8pt}{8pt}}
\makeatother

\makeatletter
\newcommand{\srcsize}{\@setfontsize{\srcsize}{6pt}{6pt}}
\makeatother

\makeatletter
\newcommand{\srcsmallsize}{\@setfontsize{\srcsmallsize}{5pt}{5pt}}
\makeatother

% to not copy line numbers
\newcommand{\noncopynumber}[1]{%
  \BeginAccSupp{method=escape,ActualText={}}%
    \color{blue} #1%
  \EndAccSupp{}%
}
% to not copy line numbers

\title[C++]
{Программирование на языке C++}
 
\subtitle{Вводный курс}
 
\author[А.~Б.~Морозов]
{
  \texorpdfstring{Александр Морозов\newline\href{mailto:gelu.speculum@gmail.com}{gelu.speculum@gmail.com}}
  {Александр Морозов}
}
  
\date[ITMO 2020]
{ИТМО, весенний семестр 2020}
 
\logo{%
  \makebox[0.97\paperwidth]{%
    \includegraphics[align=c,width=2cm,keepaspectratio]{itmo_logo.png}
    \hfill
    \includegraphics[align=c,width=1.5cm,keepaspectratio]{itiviti_logo.png}
  }
}

\AtBeginSection[]
{
  \begin{frame}
    \frametitle{Содержание}
    \tableofcontents[currentsection]
  \end{frame}
}

\begin{document}

\frame{\titlepage}

%-------------------------------------------------
\section{Пространства имён}

\begin{frame}[fragile]{Пространства имён}
  Позволяют использовать одинаковые имена для разных сущностей в разных частях программы.

  \begin{lstlisting}
  namespace A {
    void f(int);
    double g();
  }
  namespace B {
    bool f(const std::string &);
    void g();
  }
  int main()
  {
    int i = 101;
    A::f(i);
    double x = A::g();
    B::g();
  }
  \end{lstlisting}
\end{frame}

\begin{frame}[fragile]{Пространства имён и определения}
  \begin{lstlisting}
  namespace A {
    int f();
    void g();
    double sqrt(double x)
    {
      ...
    }
  }

  int A::f()
  { ... }

  void A::g()
  { ... }
  \end{lstlisting}
\end{frame}

\begin{frame}[fragile]{Вложенные пространства имён}
  \begin{lstlisting}
  namespace A {
    namespace B {
      namespace C {
        class C {...};
      }
    }
  }
  namespace A::B::C {
    struct S {...};
  }

  int main()
  {
    A::B::C::C c;
    A::B::C::S s;
  }
  \end{lstlisting}
\end{frame}

\begin{frame}[fragile]{Область видимости пространства имён}
  \begin{lstlisting}
  namespace A {
    const int n = 1;

    namespace B {
      void f()
      { std::cout << n << std::endl; }
    }
  }

  const int n = 0;

  namespace A {
    namespace B {
      const int n = 3;

      void g()
      { f(); }

      void h()
      { std::cout << n << std::endl; }
    }
  }

  int main()
  {
    A::B::g();
    A::B::h();
  }
  \end{lstlisting}
\end{frame}

\begin{frame}[fragile]{Область видимости пространства имён, продолжение}
  \begin{lstlisting}
  namespace A {
    int f(int);
    const int n = -13;
  }

  int A::f(int i = n)
  { return i; }
  \end{lstlisting}
\end{frame}

\begin{frame}[fragile]{Анонимные пространства имён}
  \begin{lstlisting}
  namespace {
    class C {...};
    C c;
    void f() {...}
  }
  namespace A {
    namespace {
      int g() {...};
    }
  }

  int main()
  {
    C cc = c;
    f();
    return A::g();
  }
  \end{lstlisting}
\end{frame}

\begin{frame}[fragile]{Полностью квалифицированные имена}
  \begin{lstlisting}
  namespace U {
    namespace X {
      void f();
    }
    namespace Y {
      namespace X {
        double f(int);
      }
      void g(int n) {
        double x = X::f(n);
        ::U::X::f();
      }
    }
  }
  \end{lstlisting}
\end{frame}

\begin{frame}[fragile]{Псевдонимы пространств имён}
  \begin{lstlisting}
  namespace A::B::C::D::E {
    class X {...};
  }
  namespace abcde = A::B::C::D::E;

  int main()
  {
    abcde::X x;
  }
  \end{lstlisting}
\end{frame}

\begin{frame}[fragile]{using директива и объявление}
  \begin{lstlisting}
  namespace A {
    using namespace std; // using directive

    void say_hi()
    {
      cout << "Hi!" << endl;
    }
  }
  namespace B {
    using std::string; // using declaration

    std::size_t strlen(const string & s)
    {
      return s.size();
    }
  }
  \end{lstlisting}
\end{frame}

\begin{frame}[fragile]{Особенности using}
  \begin{lstlisting}
  class X {};

  void swap(X &, X &) {}

  namespace N {
    void swap(X &, X &) {}
  }

  struct Y { X x; };

  void swap(Y & l, Y & r) {
    using std::swap;
    using namespace N; // N::swap is hidden
    swap(l.x, r.x); // ::swap is found by ADL
  }
  \end{lstlisting}
\end{frame}

\begin{frame}{Добавление в пространство std}
  В общем случае добавление в пространство \lstinline{std} запрещено и является UB.

  \vspace{1em}

  Исключение: специализация шаблонных классов из пространства \lstinline{std}, которая зависит хотя бы от одного пользовательского типа.

  \vspace{1em}

  Ограничения добавления специализаций:
  \begin{itemize}
    \item полная специализация метода класса из \lstinline{std} -- UB
    \item полная или частичная специализация класса, являющегося вложенным в класс из \lstinline{std} -- UB
  \end{itemize}
\end{frame}

\begin{frame}[fragile]{Пример правильного расширения std}
  \begin{lstlisting}[basicstyle=\ttfamily\srcbigsize]
  struct S
  {
    std::string str;
  };

  namespace std {
    template <>
    struct hash<S> : hash<std::string> {
      using argument_type = S;
      using result_type = size_t;
      size_t operator() (const S & s) const noexcept {
        return hash<std::string>::operator() (s.str);
      }
    };
  }

  void f() {
    std::unordered_set<S> set;
    set.insert("foo");
    set.insert("foo");
    set.insert("bar");
    assert(set.size() == 2);
  }
  \end{lstlisting}
\end{frame}

%-------------------------------------------------
\section{Поиск имён}

\begin{frame}{Виды поиска имён}
  \begin{itemize}
    \item Квалифицированный \vspace{2em}
    \item Неквалифицированный \vspace{2em}
    \item ADL
  \end{itemize}
\end{frame}

\begin{frame}{Квалифицированный поиск имён}
  Квалифицированное имя -- имя справа от оператора \lstinline{::}.

  \vspace{1em}

  Может относиться к:

  \vspace{0.5em}

  \begin{itemize}
    \item члену класса \vspace{0.5em}
    \item сущности из пространства имён \vspace{0.5em}
    \item элементу \lstinline{enum}
  \end{itemize}

  \vspace{1em}

  Слева от оператора \lstinline{::} -- либо ничего, либо уже найденное имя (класса, пространства имён или \lstinline{enum}).
\end{frame}

\begin{frame}[fragile]{Примеры квалифицированного поиска имён}
  \begin{lstlisting}
  namespace N {   
  struct A {
    int f() const { return 10; }
  };
  struct B : A {
    int f() const { return 15; }

    enum class E { A, B, C };
  };
  }

  void g() {
    N::B b;
    std::cout << b.f() << std::endl;
    std::cout << b.A::f() << std::endl;
    std::cout << static_cast<int>(N::B::E::A) << std::endl;
  }
  \end{lstlisting}
\end{frame}

\begin{frame}[fragile]{Приоритет поиска в пространствах имён}
  \begin{lstlisting}
  namespace A {
    void f();
    void g();
  }
  namespace B {
    using namespace A;
    void f();
  }

  int main() {
    B::f(); // 'f' from 'B'
    B::g(); // 'g' from 'A'
  }
  \end{lstlisting}
\end{frame}

\begin{frame}[fragile]{Elaborated type specifier}
  \begin{lstlisting}
  class X {
  public:
    struct Y {};
  private:
    int Y;
  };

  void f() {
    char X;
    X x; // error, 'X' refers to variable X
    class X xx; // OK
    X::Y y; // error, 'X::Y' refers to private member 'Y'
    class X::Y yy; // OK
  }
  \end{lstlisting}
\end{frame}

\begin{frame}{Неквалифицированный поиск}
  Если имя не стоит справа от \lstinline{::} применяется неквалифицированный поиск.

  \vspace{1em}

  При таком поиске будут анализироваться только подходящие сущности (например, для имени слева от \lstinline{::} -- только пространства имён, классы и \lstinline{enum}).

  \vspace{1em}

  Такой поиск обрабатывает \lstinline{using} директивы, как если бы содержимое соответствующего пространства имён находилось в ближайшем окружающем пространстве имён.

  \vspace{1em}

  Поиск будет идти вплоть до обнаружения хотя бы одной сущности с таким именем или до исчерпания вариантов. Результатом будут все подходящие сущности с искомым
  именем, объявленные на первом уровне, на котором обнаружилось имя.
\end{frame}

\begin{frame}[fragile]{ADL: пример}
  \begin{lstlisting}
  std::cout << x << std::endl;
  operator << (std::cout, x);
  operator << (std::cout, std::endl);

  std::cout << endl; // error: 'endl' identifier not found
  endl(std::cout); // OK
  \end{lstlisting}
\end{frame}

\begin{frame}[fragile]{ADL: пример посложнее}
  \begin{lstlisting}[basicstyle=\ttfamily\srcbigsize]
  namespace adl {
    struct X {
      int n = 10;
      X & operator + (const X & other) {
        n += other.n;
        return *this;
      }
    };
    X operator + (const X & lhs, const int rhs)
    {
      X x = lhs;
      x.n += rhs;
      return x;
    }
    enum class E { A, B, C };
    std::ostream & operator << (std::ostream & strm, const E & e)
    { return strm << static_cast<int>(e); }
  }

  int main() {
    adl::X x1, x2 = adl::X{} + 10;
    std::cout << (x1 + x2).n << std::endl;
    adl::E e = adl::E::B;
    std::cout << e << std::endl;
  }
  \end{lstlisting}
\end{frame}

\begin{frame}{ADL: правила}
  Применяется для имени из выражения вызова функции.

  \vspace{1em}

  Не производится, если неквалифицированный поиск дал:
  \begin{itemize}
    \item член класса
    \item объявление функции на уровне блока (не считая \lstinline{using} объявления)
    \item объявление сущности, не являющееся функцией
  \end{itemize}

  \vspace{1em}

  Иначе каждый аргумент функционального вызова добавляет пространства имён и/или классы во множество, в котором затем будет производиться поиск (помимо обычного неквалифицированного поиска).

  \vspace{1em}

  В процессе поиска отбираются только объявления функций.
\end{frame}

\begin{frame}{ADL: правила для различных типов аргументов}

  В зависимости от типа аргумента к поиска добавляется:

  \begin{itemize}
    \item базовый тип -- ничего
    \item класс -- сам класс, его предки, его область определения (окружающий класс или пространство имён)
    \item указатель на функцию -- анализ применяется для типа возвращаемого значения и типов всех параметров
    \item шаблон -- в дополнение к обычным правилам, анализ применяется для каждого шаблонного параметра-типа
    \item \lstinline{enum} -- область определения \lstinline{enum}
    \item \lstinline{T *} -- анализ применяется для \lstinline{T}
    \item указатель на метод -- анализ применяется для класса, типа возвращаемого значения и всех параметров
    \item указатель на поле -- анализ применяется для типа поля и класса
  \end{itemize}
\end{frame}

\end{document}
